%% 06_conclusion.tex — Conclusión y trabajo futuro.

\section{Conclusión y trabajo futuro}
\label{sec:conclusion}

Se ha presentado \herramienta, un motor DAST asíncrono cuya arquitectura separa
explícitamente la planificación de \emph{payloads} y la confirmación de
hallazgos en etapas desacoplables, junto con una metodología de evaluación
reproducible sobre OWASP Benchmark, un oráculo independiente, un diseño RCBD y
un manifiesto con \emph{hash} del corpus.

El estudio de ablación completo —\NumCorridas{}~ejecuciones sobre la rejilla de
cuatro presupuestos y \NumCondiciones{}~configuraciones— arroja un resultado
honesto y, en última instancia, más útil que el que anticipaba nuestra
hipótesis. La mayoría de las ablaciones son \emph{nulas}: el entrelazado por
contexto, la confirmación en tiempo de ejecución y la ordenación adaptativa en
vivo no modifican la exhaustividad, la precisión ni la FPR en este banco; la
prueba de Friedman detecta diferencias en el \emph{orden} de las peticiones
($p = \ValorFriedmanP{}$), pero con tamaños del efecto insignificantes
($\lvert\delta\rvert < \num{0.02}$).

El valor del trabajo se desplaza entonces hacia la \textbf{economía de
recursos} del escaneo (\cref{sec:res-sintesis}): el $F_1$ es máximo en un
presupuesto ajustado de \num{20}~peticiones por parámetro (\ValorFOneBVeinte{},
exhaustividad \ValorRecallTecho{}, FPR \ValorFPRBaseline{}); ampliarlo a
\num{50} o más no detecta \emph{ninguna} vulnerabilidad adicional pero suma
\SaturacionFPnuevos{}~falsos positivos ($-\DeltaPrecisionSaturacion{}$ de
precisión); y la priorización por confianza —única etapa con efecto práctico—
solo preserva la selectividad en ese régimen de escasez
($-\DeltaPrecisionNoPriority{}$ de precisión si se desactiva a presupuesto
\num{20}, nada a partir de \num{50}). La conclusión operativa es que optimizar
un DAST moderno no consiste en disparar \emph{payloads} sin límite, sino en
\textbf{maximizar la eficiencia bajo escasez}: fijar un presupuesto ajustado y
gastarlo en el orden correcto. La confirmación en tiempo de ejecución no
penaliza la precisión porque en este banco los indicadores crudos ya son
fiables.

Frente a todo esto, la barrera dura es el \textbf{techo de recuperación}: los
verdaderos positivos se estancan en \ValorTPtecho{}$/$\NumVulnerables{} y los
\ValorFNtecho{}~falsos negativos persisten a presupuesto $\infty$. No es un
problema de planificación —el planificador solo ordena un corpus que a
presupuesto ilimitado se emite entero— sino de \textbf{cobertura de vectores}:
para esas instancias el corpus no contiene ninguna carga que dispare la
condición vulnerable bajo su contexto de salida. Romper ese techo exige
ampliar la biblioteca de \emph{payloads}, no refinar las heurísticas de
ordenación.

\noindent\textbf{Trabajo futuro.}
\begin{enumerate}
  \item \textbf{Expandir la biblioteca de \emph{payloads}} para atacar las
  \ValorFNtecho{}~instancias fuera de alcance —caracterizando sus contextos de
  salida y codificaciones e incorporando vectores dependientes de codificación
  y \emph{polyglots}—: es la vía con mayor rendimiento marginal esperado, por
  encima de seguir afinando la ordenación.
  \item Repetir el estudio sobre aplicaciones reales y bancos con señales
  ruidosas —donde la confirmación en tiempo de ejecución debería reducir
  falsos positivos que aquí no fueron observables— y sobre vectores
  \code{POST}, cabeceras y \emph{cookies}.
  \item Refinar la rejilla alrededor del punto óptimo (presupuestos
  $\in [\num{10}, \num{30}]$) con réplicas, para localizar el máximo de $F_1$
  con más precisión, y añadir un precalentamiento explícito de la JVM para el
  análisis de latencia absoluta.
\end{enumerate}
